\documentclass[UTF8]{ctexart}
\usepackage{xeCJK}
\usepackage{geometry}
\geometry{left=2cm,right=2cm,top=2.5cm,bottom=2.5cm}
\setlength{\headheight}{12.7pt}

\usepackage{titlesec}
\titleformat{\section}
  {\normalfont\large\bfseries}
  {\thesection}
  {1em}
  {}
\titlespacing*{\section}
  {0pt}
  {3.5ex plus 1ex minus .2ex}
  {2.3ex plus .2ex}

\usepackage{amsmath}
\usepackage{amssymb}
\usepackage{amsthm}
\usepackage{enumitem}
\setlist[enumerate]{itemsep=0pt,topsep=0pt,parsep=0pt,leftmargin=0pt,itemindent=2.5em}
\setlist[itemize]{itemsep=0pt,topsep=0pt,parsep=0pt,leftmargin=0pt,itemindent=2.5em}
\usepackage{fancyhdr}
\pagestyle{fancy}
\fancyhf{}
\usepackage{graphicx}
\usepackage{hyperref}
\usepackage{indentfirst}
\usepackage{lmodern}


\begin{document}
\setlength{\abovedisplayskip}{1pt}
\setlength{\belowdisplayskip}{1pt}

\section{群作用}

% 收到了博士生录取通知书, 大概是此生的最后一张录取通知书.
% 很巧的是, 一年前恰好刚刚结束推研面试. 时光飞逝啊.
% 2026.07.18

\hypertarget{sec:定义1.1}{}
\noindent\textbf{定义1.1 (群作用)}

给定群$G$和集合$X$, 以及映射$\cdot:G\times X\to X$. 如果
\begin{enumerate}
    \item 对任意的$x\in X$, 有$1_G\,x=x$,
    \item 对任意的$g_1,g_2\in G$和$x\in X$, 有$g_1(g_2x)=(g_1g_2)x$,
\end{enumerate}
就称$\cdot$是$G$在$X$上的一个\textbf{群作用}, 记作$\cdot:G\curvearrowright X$.

% 只睡了两个多小时...
% 恭喜西班牙! 早安.
% 2026.07.20

\vspace{10pt}

如果此时有$H<G$, 显然$\cdot\!\restriction_{H\times X}$可以给出$H$在$X$上的群作用.

\vspace{10pt}

\hypertarget{sec:定理1.1}{}
\noindent\textbf{定理1.1}

任给群$G$和集合$X$, 对任意的$\cdot:G\curvearrowright X$, 有群同态
\[
    \overline{\cdot}:G\to S_X,\quad g\mapsto(x\mapsto gx),
\]
并且$\cdot\mapsto\overline{\cdot}$是
$G$在$X$上的群作用与$G$到$S_X$的群同态之间的一一对应.

\vspace{10pt}

\hypertarget{sec:定义1.2}{}
\noindent\textbf{定义1.2}

给定群作用$\cdot:G\curvearrowright X$, 如果$\overline{\cdot}$是单同态, 即
\[
    \{g\in G\mid\forall x\in X:gx=x\}=\mathrm{Ker}\,\overline{\cdot}=\{1_G\},
\]
那么称$\cdot$是\textbf{忠实}的.

\vspace{10pt}

对于一般的群作用$\cdot$\,, $\mathrm{Ker}\,\overline{\cdot}$表示群中
不产生任何作用的元素, 其商群就可以产生忠实的群作用.

\vspace{10pt}

\hypertarget{sec:定理1.2}{}
\noindent\textbf{定理1.2}

对任意的群作用$\cdot:G\curvearrowright X$, 由商群的泛性质, 有群同态
\vspace{3pt}
\[
    \overline{\times}:\,\!^{\displaystyle G}\!
    /_{\displaystyle\mathrm{Ker}\,\overline{\cdot}}\,\to S_X,\\[3pt]
\]
它对应了忠实的群作用$\times:\,\!^{\displaystyle G}\!
/_{\displaystyle\mathrm{Ker}\,\overline{\cdot}}\curvearrowright X$, 并且
对任意的$g\in G$和$x\in X$有$g\mathrm{Ker}\,\overline{\cdot}\times x=gx$.

\vspace{3pt}
\noindent{\kaishu 证明.}

由商群的泛性质, $\mathrm{Ker}\,\overline{\times}=
\,\!^{\displaystyle\mathrm{Ker}\,\overline{\cdot}}\!/_
{\displaystyle\mathrm{Ker}\,\overline{\cdot}}=\{\mathrm{Ker}\,\overline{\cdot}\}$,
从而$\times$忠实, 并且对任意的$g\in G$和$x\in X$有
\vspace{5pt}
\begin{equation*}
    g\mathrm{Ker}\,\overline{\cdot}\times x=
    \overline{\times}(g\mathrm{Ker}\,\overline{\cdot})(x)=
    \overline{\cdot}(g)(x)=gx.
\tag*{\qedsymbol}\end{equation*}

\vspace{15pt}

下面举一些常见的群作用的例子.

\vspace{10pt}

\hypertarget{sec:例1.1}{}
\noindent\textbf{例1.1}

任给群$G$, 它在$G$自身上有群作用$(g,g')\mapsto gg'$, 称之为左乘作用.

\vspace{10pt}

\hypertarget{sec:例1.2}{}
\noindent\textbf{例1.2}

任给群$G$和它的正规子群$H$, $G$在$H$上有群作用$(g,h)\mapsto ghg^{-1}$, 称之为共轭作用.

\vspace{10pt}

\hypertarget{sec:例1.3}{}
\noindent\textbf{例1.3}

任给群$G$和它的子群$H$, $G$在$\,\!^{\displaystyle G}\!/_{\displaystyle H}$
上有群作用$(g,P)\mapsto gP$, 称之为对陪集的左乘作用.





\newpage

\section{\texorpdfstring{$G$}{G}-集与等变映射}

\hypertarget{sec:定义2.1}{}
\noindent\textbf{定义2.1 ($G$-集范畴)}

给定群$G$. 对任意的集合$X$和其上的群作用$\cdot:G\curvearrowright X$,
称$(X,\cdot)$或$X$是一个$G$-\textbf{集}.

给定$G$-集$(X,\cdot_X)$, $(Y,\cdot_Y)$和映射$f:X\to Y$. 如果
\[
    \forall g\in G,\,\forall x\in X:\quad f(g\cdot_X x)=g\cdot_Y f(x),
\]
就称$f$是从$(X,\cdot_X)$到$(Y,\cdot_Y)$的\textbf{等变映射}.

若干$G$-集与它们之间的等变映射可以组成范畴, 称作$G$-\textbf{集范畴}.

% 忽然意识到, 这是线性空间和线性映射的最简单版本.

\vspace{10pt}

\hypertarget{sec:定理2.1}{}
\noindent\textbf{定理2.1}

任给群$G$, $G$-集$(X,\cdot_X)$, $(Y,\cdot_Y)$和它们之间的等变映射$f$.
如果$f$是双射, 那么$f^{-1}$也是等变映射.

\noindent{\kaishu 证明.}

对任意的$g\in G$和$y\in Y$, 记$x=f^{-1}(y)$, 则有
\begin{equation*}
    g\cdot_Y y=g\cdot_Y f(x)=f(g\cdot_X x)\quad\implies\quad
    f^{-1}(g\cdot_Y y)=g\cdot_X x=g\cdot_X f^{-1}(y).
\tag*{\qedsymbol}\end{equation*}

\vspace{10pt}

据上述定理, 如果两个$G$-集之间有等变的双射, 那么它们在$G$-集范畴中同构.

\vspace{15pt}

\hypertarget{sec:定义2.2}{}
\noindent\textbf{定义2.2 ($G$-子集)}

给定群$G$, $G$-集$X$和$X'\subset X$, 如果
\[
    \forall g\in G,\,\forall x\in X':\quad gx\in X',
\]
那么显然$X'$也是$G$-集, 称之为$X$的$G$-\textbf{子集}.
显然这时$X'$到$X$的嵌入映射是等变映射.

\vspace{10pt}

\hypertarget{sec:定理2.2}{}
\noindent\textbf{定理2.2}

任给群$G$, $G$-集$X,Y$和它们之间的等变映射$f$, 则
\begin{enumerate}
    \item 对$X$的任意$G$-子集$X'$, $f[X']$是$Y$的$G$-子集,
    \item 对$Y$的任意$G$-子集$Y'$, $f^{-1}[Y']$是$X$的$G$-子集.
\end{enumerate}

\noindent{\kaishu 证明.}

\begin{enumerate}
    \item 对任意的$g\in G$和$y\in f[X']$, 存在$x\in X'$使得$y=f(x)$, 从而
    $gy=gf(x)=f(gx)\in f[X']$,
    \item 对任意的$g\in G$和$x\in f^{-1}[Y']$, 有$f(x)\in Y'$, 从而
    $f(gx)=gf(x)\in Y'$, 即$gx\in f^{-1}[Y']$. \hfill $\qedsymbol$
\end{enumerate}

\vspace{10pt}

\hypertarget{sec:定义2.3}{}
\noindent\textbf{定义2.3 (生成$G$-子集)}

任给群$G$, $G$-集$X$和$S\subset X$, 定义它的\textbf{生成}$G$-\textbf{子集}为
\[
    \langle S\rangle\overset{\mathrm{def}}{=}\{gs\mid g\in G,s\in S\},
\]
它是$X$的包含$S$的最小$G$-子集.

\noindent{\kaishu 验证.}

一方面, 对任意的$h\in G$和$gs\in\langle S\rangle$, 都有
$h(gs)=(hg)s\in\langle S\rangle$, 所以$S$是$G$-子集且显然包含$S$.

另一方面, 如果$X'$是$X$的包含$S$的$G$-子集, 那么对任意的$g\in G$和$s\in S$都有
\[
    s\in X'\quad\implies\quad gs\in X',
\]
从而$\langle S\rangle\subset X'$. \hfill $\qedsymbol$

\end{document}